\documentclass[tikz,border=10pt]{standalone}
\input{../../../docs/tikz/dag_preamble.tex}

\begin{document}
\begin{tikzpicture}[
  x=1cm,
  y=1cm,
  every node/.append style={outer sep=4pt}
]

\dagPaperMetadata
  {A No Free Lunch Theorem for Human-AI Collaboration}
  {Peng, Garg, Kleinberg}
  {AAAI}
  {2025}
  {https://ojs.aaai.org/index.php/AAAI/article/view/33574}

\dagPaperLegendRightOfMetadata{
\node[dag_result, dag_template_legend] (legRes) at (0,0) {checked named\\result};
\node[dag_lemma, dag_template_legend] (legLem) at (4.0,0) {checked\\construction};
\node[dag_model, dag_template_legend] (legDef) at (8.0,0) {checked finite\\definition};
\node[dag_partial, dag_template_legend] (legPart) at (12.0,0) {partial source\\coverage};
\node[dag_template_note, dag_template_legend, text width=2.72cm]
  (legNote) at (16.0,0) {formalized source note\\(not a caveat)};
}
\daglegend{(legRes)(legLem)(legDef)(legPart)(legNote)}{Legend}

\begin{dagPaperBody}

\node[font=\sffamily\bfseries\Large, align=center, text width=22cm]
  (Class) at (6,0.6) {
  Status: partially formalized because of three coverage boundaries;\\
  Theorem 1 is checked and no central-result caveat is claimed
};

% The three yellow nodes are the only drivers of partial status.
\node[dag_partial, text width=5.2cm] (GeneralDomain) at (-5,-2.0) {
  \textbf{Partial boundary 1} \\
  full arbitrary probability-space\\
  collaboration-setting universe\\
  \emph{abstract inclusion bridge checked}
};

\node[dag_model, text width=5.2cm] (Finite) at (2,-2.0) {
  \textbf{Finite calibrated settings} \\
  mass, label probability, predictions,\\
  calibration, induced accuracies
};

\node[dag_partial, text width=5.2cm] (Partitions) at (9,-2.0) {
  \textbf{Partial boundary 2} \\
  general partition-to-predictor recipe\\
  \emph{every proof partition checked}
};

\node[dag_partial, text width=5.2cm] (Epsilon) at (16,-2.0) {
  \textbf{Partial boundary 3} \\
  full $\epsilon<1/2$ S1 family\\
  \emph{$\epsilon=1/4$ checked}
};

\node[dag_lemma, text width=4.8cm] (Mixture) at (-5,-6.4) {
  \textbf{Proposition 6} \\
  disjoint-union mixture; mass, label,\\
  predictor, and accuracy equations
};

\node[dag_lemma, text width=4.8cm] (Lemma8) at (2,-6.4) {
  \textbf{Lemma 8} \\
  odds-mass witness, calibration,\\
  weak/strict accuracy dominance
};

\node[dag_lemma, text width=4.8cm] (S1) at (9,-6.4) {
  \textbf{Proposition 9: $S_1$} \\
  checked $\epsilon=1/4$ instance;\\
  accuracies $7/12$ and $3/4$
};

\node[dag_lemma, text width=4.8cm] (S2) at (16,-6.4) {
  \textbf{Proposition 9: $S_2$} \\
  paired odds witnesses, calibration,\\
  and strict gap for agent $k$
};

\node[dag_result, text width=5.3cm] (Prop7) at (-0.5,-10.8) {
  \textbf{Proposition 7} \\
  finite reliability forces one fixed\\
  off-half deferral coordinate
};

\node[dag_result, text width=5.3cm] (Prop9) at (12.5,-10.8) {
  \textbf{Proposition 9} \\
  finite reliability forces one fixed\\
  tie label; final weights $7/8,1/8$
};

\node[dag_result, text width=8.5cm] (Main) at (6,-15.0) {
  \textbf{Theorem 1 --- checked forward implication} \\
  source reliability in any domain containing all finite\\
  calibrated witnesses implies non-collaboration
};

\draw[dag_dashed_arrow] (GeneralDomain.east) -- node[above, font=\scriptsize]
  {accuracy-preserving inclusion} (Finite.west);
\draw[dag_arrow] (Finite.south west) -- (Mixture.north east);
\draw[dag_arrow] (Finite.south) -- (Lemma8.north);
\draw[dag_dashed_arrow] (Partitions.south west) -- (Lemma8.north east);
\draw[dag_arrow] (Finite.south east) -- (S1.north west);
\draw[dag_arrow] (Finite.east) -- ++(1.4,0) |- (S2.north);
\draw[dag_dashed_arrow] (Partitions.south east) -- (S2.north west);
\draw[dag_dashed_arrow] (Epsilon.south west) -- (S1.north east);
\draw[dag_arrow] (Mixture.south) -- ++(0,-1.0) -| (Prop7.north west);
\draw[dag_arrow] (Lemma8.south) -- (Prop7.north east);
\draw[dag_arrow] (Mixture.south east) -- ++(0,-0.8) -| (Prop9.north west);
\draw[dag_arrow] (S1.south) -- (Prop9.north west);
\draw[dag_arrow] (S2.south) -- (Prop9.north east);
\draw[dag_arrow] (Prop7.south) -- ++(0,-1.0) -| (Main.north west);
\draw[dag_arrow] (Prop9.south) -- ++(0,-1.0) -| (Main.north east);

% Gray nodes record source clarifications and never drive partial/caveat status.
\node[font=\sffamily\bfseries\large, align=center] at (6,-18.0)
  {Three source clarifications --- not caveats and not partial-status drivers};

\node[dag_template_note, text width=6.1cm] (LossNote) at (-4.5,-20.0) {
  \textbf{Note 1: loss/accuracy label.}
  $\mathrm{E}|\widehat Y-Y|$ is loss; Lean proves expected correctness is its
  complement.
};

\node[dag_template_note, text width=6.1cm] (TieNote) at (6,-20.0) {
  \textbf{Note 2: half-tie strictness.}
  The strict-gap form is used for $\eta\ne1/2$; the $\eta=1/2$ boundary is
  treated as a tie.
};

\node[dag_template_note, text width=6.1cm] (IffNote) at (16.5,-20.0) {
  \textbf{Note 3: Definition 4 boundary.}
  Definition 4 exempts boundary profiles; the displayed Theorem 1 forward
  implication is unchanged.
};

\node[dag_template_note, text width=20.5cm, align=center] (Inventory) at (6,-23.3) {
  \textbf{Selected source accounting: 45 atomic items.}
  36 direct-routed items and 6 support-only items; the remaining source notes
  are recorded as scoped mathematical clarifications.
};

\end{dagPaperBody}
\end{tikzpicture}
\end{document}
